\begin{table}[ht]
\caption{\textbf{\sustainable example}. The positive document explains the critical concept of soluble salts, which indicate that they contain dissolved minerals and could be harmful to plants.}
\centering
% \resizebox{\textwidth}{!}{
\begin{tabular}{p{13cm}}
\toprule
% \hspace{5.5cm} \sustainable \\
% \midrule
\textit{\textbf{Query}} \\
\midrule
How good is it to reuse water from plant pots?

I'm living in an apartment, and after I water my plants the water goes to plates below the pots. The pots are in a metallic structure above the plates, so I can take the plates to reuse the water (throwing it at the plants again).

This reuse seems beneficial, because I think I can get rid of mosquitoes that would reproduce in the stagnated water. And also some nutrients of the soil (as well as earthworms) can return to the vase.

Is there some negative points in doing that? \\
\midrule
\textit{\textbf{Chain-of-thought reasoning to find documents}} \\
\midrule
The water that accumulates in the plates as a result of watering is likely to contain minerals, soluble salts, and other materials that exist in soil or plants. 
To figure out whether there is any negative point in reusing that water, we thus need to understand whether the components in that water will result in any adverse effects. \\
\midrule
\textit{\textbf{Example positive document}} \\
\midrule
Soluble Salts

Soluble salts may accumulate on the top of the soil, forming a yellow or white crust. A ring of salt deposits may form around the pot at the soil line or around the drainage hole. Salts may also build up on the outside of clay pots. In house plants, signs of excess soluble salts include reduced growth, brown leaf tips, dropping of lower leaves, small new growth, dead root tips, and wilting.

Soluble salts are minerals dissolved in water. Fertilizer dissolved in water becomes a soluble salt. When water evaporates from the soil, the minerals or salts stay behind. As the salts in the soil become more and more concentrated, it becomes more difficult for plants to take up water. If salts build up to an extremely high level, water can be taken out of the root tips, causing them to die. High levels of soluble salts damage the roots directly, weakening the plant and making it more susceptible to attack from insects and diseases. One of the most common problems associated with high salt levels is root rot.

The best way to prevent soluble salt injury is to stop the salts from building up. When watering, allow some water to drain through the container and then empty the saucer. Do not allow the pot to sit in water. If the drained water is absorbed by the soil, the salts that were washed out are reabsorbed through the drainage hole or directly through a clay pot. 

...\\

% House plants should be leached at least every 4 to 6 months. To leach plants, pour excess water on the soil and let it drain completely. The amount of water used for leaching should equal twice the volume of the pot. Keep the water running through the soil to wash out the salts. If a layer of salts has formed a crust at the soil surface, remove the salt crust before leaching. Do not remove more than 1/4 inch (6 mm) of soil. It is best not to add soil to the top of the pot. If the soluble salt level appears to be extremely high, repot the plant.

% The level of salts that causes injury varies with the species of plant and how it is grown. A house plant may be injured by salts at a very low concentration, but the same plant growing in a greenhouse where watering is well managed may tolerate salts at high levels. Some nurseries and retail plant outlets leach plants to remove excess salts before the plant is sold. If you are not sure that a newly purchased plant has been leached, leach it the first time you water. \\
\midrule
\textit{\textbf{Example negative document}} \\
\midrule
Water reuse in California

Water reuse in California is the use of reclaimed water for beneficial use. As a heavily populated state in the drought-prone arid west, water reuse is developing as an integral part of water in California enabling both the economy and population to grow.

Wastewater
Reclaimed water is treated wastewater that comes from homes and businesses, such as sink water, shower water, and toilet water including everything dumped into wastewater drains from laundry soap to bleach to oil to human waste. Wastewater can be divided into greywater and blackwater, with the first being defined as water that had been used for laundry, bathing, sink washing, and dishwashers. Blackwater is defined as sewage that includes feces from toilets.[1] Due to the low amounts of physical pollutants in greywater, most of its contaminants are dissolved organic matter, which can be physically filtered and cleaned through various membranes, as well as through biological treatment methods. 

...\\

% The subsequent heterogeneous solution is collected through pipes and the sewer system, and is then treated at a wastewater treatment plant to the standard required according to its intended use. Historically, this recycled water has been used for agriculture, large-scale landscaping, industrial processes like cooling systems, and groundwater recharge. Future infrastructure for water reuse may include washing laundry and cars or watering lawns and flushing our toilets, in addition to use in municipal infrastructure for street cleaning, fountains, and commercial use.[3] Water in California is to be used 'reasonably' under the Reasonable Use doctrine, and not reusing water when possible constitutes a violation of this doctrine pursuant to Water Code sections 13550 et seq. according to the State Water Resources Control Board.[4] Water reuse in California has more than tripled since the 1970s, growing from less than 200,000 acre-feet per year to nearly 700,000 in 2009.
\bottomrule
\end{tabular}
% }
\label{tab:example_sustainable_living}
\end{table}